\chapter{School Boards and the Democratic Ideal}
\label{chap:conclusion}
\chapterstyle{deposit}

\section{Introduction}

School boards are a puzzling feature of the American polity worth considering. 
Collectively, school boards manage annual public school expenditures rivaling the 
budget of the Department of Defense. School boards set policy and guidance on a 
host of issues that affect the daily lives of nearly fifty million school children 
in public schools and their parents. School boards make decisions about the 
books that are taught, the graduation requirements that must be met, and the 
distribution of resources within and between schools in a community. All of this, 
and more, is entrusted to tens of thousands of local elected officials in school 
districts across the country. 

School boards are one of the most high-profile examples of America's most local 
form of governments -- special purpose districts. Special purpose districts are 
a local government with authority over a specific issue such as water rights, 
transportation policy, or education. Often, special purpose districts have their 
own boundaries that are not coterminous with county or municipal government 
boundaries, and often they are governed by non-partisan elected officials. While 
responsible for a large share of the total population of elected officials in the 
the United States, special purpose governments, including school boards, have been 
largely unstudied by scholars of American politics. 

This dissertation represents a first step toward a more comprehensive picture of 
school boards as democratic institutions. 

\subsection{Key Findings}

Using Wisconsin school board elections from the 2002-2012 period, I have shown 
that school board seats are unlikely to be heavily contested with few school 
districts ever holding a primary. Fewer than half of all elections featured even 
a nominal challenge or a ballot where voters had a non-write in option. 

Incumbents were featured in 83\% of races, and in 48\% of all races, voters 
had a choice between an incumbent and a challenger. As a comparison, these figures 
represent a level of contestation similar to partisan primary elections for the 
U.S. House. 

Despite this, if a school board race is contested, the contest is often quite 
close. This has more to do with the small electorate than the level of campaign 
activity and contestation. 

Incumbent defeat is rare, fewer than 30\% of districts had an incumbent defeated 
in any given year. Turnout is generally below 25\%. 

% Chapter 3
%


School boards do not meet a stylized standard of democracy because:

Few candidates run for office and many races are uncontested.
When incumbents face a challenger, they most often win.
When incumbents are defeated, they are defeated by a very small slice of the electorate.
But, does this make school boards truly undemocratic?

\subsection{Limitations and Future Work}

No election activity. No measure of preferences. No measure of polarization or 
division within school board preferences. 


\subsection{Future Work}

Better understand superintendent turnover and code retirements against buyouts. 

Look at financial outcomes. 

Non-electoral forms of influence. 

\subsection{Final Thought}

Theory suggests if boards are responsive to community demands and if democratic potential exists, then democratic control can be retained.

